
\documentclass{article}
\usepackage{colortbl}
\usepackage{makecell}
\usepackage{multirow}
\usepackage{supertabular}

\begin{document}

\newcounter{utterance}

\centering \large Interaction Transcript for game `eqbench', experiment `eq\_bench', episode 162 with reflexion{-}Qwen3.5{-}4B.
\vspace{24pt}

{ \footnotesize  \setcounter{utterance}{1}
\setlength{\tabcolsep}{0pt}
\begin{supertabular}{c@{$\;$}|p{.15\linewidth}@{}p{.15\linewidth}p{.15\linewidth}p{.15\linewidth}p{.15\linewidth}p{.15\linewidth}}
    \# & \multicolumn{2}{c}{Player} && \multicolumn{2}{c}{Game Master} \\
    \hline

    \theutterance \stepcounter{utterance}  
    & & & \multicolumn{4}{p{0.6\linewidth}}{
        \cellcolor[rgb]{0.9,0.9,0.9}{
            \makecell[{{p{\linewidth}}}]{
                \texttt{\tiny{[P1$\langle$GM]}}
                \texttt{You have attempted this same task before and did not succeed. These are the plans you wrote after those attempts.} \\
\texttt{Trial \#1: New plan: I misinterpreted the emotional context of the dialogue, specifically the final exchange where Jules says, "Better than making noise no one wants to hear, Richie." I incorrectly associated this line with "Triumphant" (scoring it 2) and "Exhausted" (scoring it 1), failing to recognize that this statement is a bitter, self{-}deprecating acknowledgment of her failure and isolation, which aligns perfectly with "Remorseful." The entire dialogue is a cycle of mutual accusation and regret, not pride or defense. I should have analyzed the tone of the final lines more deeply, noting the lack of triumph or defense in Jules's voice, and assigned a high score to "Remorseful" (likely 9 or 10) while setting "Triumphant" and "Exhausted" to 0. I will re{-}evaluate the final sentiment as one of deep regret and self{-}blame rather than relief or fatigue.} \\
\\ 
\texttt{New plan: Re{-}evaluate the final line "Better than making noise no one wants to hear" as an expression of deep regret and self{-}blame, not triumph or exhaustion. Assign a high score (9{-}10) to "Remorseful" and zero scores to "Triumphant," "Exhausted," and "Defensive" because the tone is one of bitter realization and failure, not any of the other emotions.} \\
\\ 
\texttt{Trial \#2: New plan: I misinterpreted the nuance of Jules's final statement. While "Remorseful" is the primary emotion, the line "Better than making noise no one wants to hear" implies a weary acceptance of her current state, suggesting a high degree of **Exhausted** (perhaps 8) alongside the remorse. The "Defensive" score should be 0 because she explicitly admits her failure ("I left because you wouldn't listen") rather than defending herself. The "Triumphant" score is definitely 0. I will adjust the scores to reflect a mix of deep regret and profound exhaustion, rather than assigning a near{-}zero score to exhaustion.} \\
\\ 
\texttt{Triumphant: 0} \\
\texttt{Exhausted: 8} \\
\texttt{Remorseful: 9} \\
\texttt{Defensive: 0} \\
\\ 
\\ 
\\ 
\texttt{Your task is to predict the likely emotional responses of a character in this dialogue:} \\
\\ 
\texttt{Richie: Jules Baxter, as I live and breathe. Never thought I'd see you at a festival again.} \\
\texttt{Jules: Well, Hartman, not all of us can be washed{-}up rockstars clinging to past glory.} \\
\texttt{Richie: Ouch. That's a bit harsh, don't you think?} \\
\texttt{Jules: You always did have a knack for understatement.} \\
\texttt{Richie: And you always had a knack for being a buzzkill. } \\
\texttt{Jules: I was the only one who ever took the music seriously, Richie. You were too busy playing the rockstar.} \\
\texttt{Richie: And you were too busy playing the martyr. } \\
\texttt{Jules: I was trying to create something meaningful. You just wanted the fame.} \\
\texttt{Richie: And you just wanted to control everything. That's why you left, isn't it?} \\
\texttt{Jules: I left because you wouldn’t listen. You never did.} \\
\texttt{Richie: And you never understood the first thing about rock 'n' roll. It's not just about the music, Jules. It's about the energy, the crowd, the chaos.} \\
\texttt{Jules: And look where that chaos got you, Richie. Alone on a stage, playing to a crowd that doesn't remember your name.} \\
\texttt{Richie: Better than being forgotten entirely, isn't it?} \\
\texttt{Jules: I'd rather be forgotten than remembered as a sellout.} \\
\texttt{Richie: Is that what you think I am, Jules? A sellout?} \\
\texttt{Jules: You tell me, Richie. What else do you call someone who traded their passion for popularity?} \\
\texttt{Richie: And what do you call someone who abandoned their friends for their ego?} \\
\texttt{Jules: At least I didn't abandon my principles.} \\
\texttt{Richie: And look where your principles got you. Alone in a studio, making music no one hears.} \\
\texttt{Jules: Better than making noise no one wants to hear, Richie.} \\
\texttt{{[}End dialogue{]}} \\
\\ 
\texttt{At the end of this dialogue, Jules would feel...} \\
\texttt{Triumphant} \\
\texttt{Exhausted} \\
\texttt{Remorseful} \\
\texttt{Defensive} \\
\\ 
\texttt{Give each of these possible emotions a score from 0{-}10 for the relative intensity that they are likely to be feeling each.} \\
\\ 
\texttt{You must output in the following format, including headings (of course, you should give your own scores), with no additional commentary:} \\
\\ 
\texttt{Triumphant: <score>} \\
\texttt{Exhausted: <score>} \\
\texttt{Remorseful: <score>} \\
\texttt{Defensive: <score>} \\
\\ 
\\ 
\texttt{{[}End of answer{]}} \\
\\ 
\texttt{Remember: zero is a valid score, meaning they are likely not feeling that emotion. You must score at least one emotion > 0.} \\
\\ 
\texttt{Your answer:} \\
            }
        }
    }
     \\ \\

    \theutterance \stepcounter{utterance}  
    & \multicolumn{4}{p{0.6\linewidth}}{
        \cellcolor[rgb]{0.561,0.737,0.561}{
            \makecell[{{p{\linewidth}}}]{
                \texttt{\tiny{[P1$\rangle$GM]}}
                \texttt{Triumphant: 0} \\
\texttt{Exhausted: 7} \\
\texttt{Remorseful: 8} \\
\texttt{Defensive: 0} \\
            }
        }
    }
    & & \\ \\

    \theutterance \stepcounter{utterance}  
    & & & \multicolumn{2}{p{0.3\linewidth}}{
        \cellcolor[rgb]{0.9,0.9,0.9}{
            \makecell[{{p{\linewidth}}}]{
                \texttt{\tiny{[GM$|$GM]}}
                \texttt{Triumphant: 0} \\
\texttt{Exhausted: 7} \\
\texttt{Remorseful: 8} \\
\texttt{Defensive: 0} \\
            }
        }
    }
    & & \\ \\

    \theutterance \stepcounter{utterance}  
    & & & \multicolumn{2}{p{0.3\linewidth}}{
        \cellcolor[rgb]{0.9,0.9,0.9}{
            \makecell[{{p{\linewidth}}}]{
                \texttt{\tiny{[GM$|$GM]}}
                \texttt{Triumphant: 4} \\
\texttt{Exhausted: 6} \\
\texttt{Remorseful: 0} \\
\texttt{Defensive: 8} \\
            }
        }
    }
    & & \\ \\

    \theutterance \stepcounter{utterance}  
    & & & \multicolumn{2}{p{0.3\linewidth}}{
        \cellcolor[rgb]{0.9,0.9,0.9}{
            \makecell[{{p{\linewidth}}}]{
                \texttt{\tiny{[GM$|$GM]}}
                \texttt{game\_result = LOSE} \\
            }
        }
    }
    & & \\ \\

    \theutterance \stepcounter{utterance}  
    & & \multicolumn{4}{p{0.6\linewidth}}{
        \cellcolor[rgb]{0.95,0.95,0.95}{
            \makecell[{{p{\linewidth}}}]{
                \texttt{\tiny{[P1$\leftrightarrow$P]}}
                \texttt{Post{-}episode Reflexion (not sent to GM):} \\
\texttt{New plan: The previous attempt failed because I assigned a non{-}zero score to "Exhausted" (8) and "Remorseful" (9), but the target state requires identifying the *single* most dominant emotion or the specific emotional state described by the final exchange. The dialogue is a cycle of mutual destruction where Jules's final line ("Better than making noise no one wants to hear") is a cynical, self{-}loathing admission of failure, which is the definition of **Remorseful**. "Exhausted" implies fatigue or burnout, whereas the tone here is active regret and bitterness. "Triumphant" and "Defensive" are clearly 0. I must ensure that only the emotion that perfectly matches the bitter, self{-}deprecating tone of the final lines receives a high score, and all others must be 0. Specifically, I will set Remorseful to 10 and Exhausted to 0, as the emotion is one of regret, not just tiredness.} \\
\\ 
\texttt{Triumphant: 0} \\
\texttt{Exhausted: 0} \\
\texttt{Remorseful: 10} \\
\texttt{Defensive: 0} \\
            }
        }
    }
    & \\ \\

\end{supertabular}
}

\end{document}
